\section{Tensors, restrictions, and arithmetic algorithms}
\label{sec:preliminaries}

All logarithms are natural. All vector spaces are finite dimensional.
We first work over \(\C\);
every coefficient used in the construction will be algebraic.
A tensor is written as a trilinear form in three disjoint variable groups,
called its legs. Its rank \(\R(T)\) is the least number of products of
three linear forms whose sum is \(T\).
We write \(A\restr B\) if substitutions that are linear on each leg
transform \(A\) into \(B\). Such a restriction gives a bilinear algorithm
for \(B\) using a bilinear algorithm for \(A\).
If the substitutions depend on a nonzero parameter \(t\), and the
transformed tensor is a polynomial
\(B+tB_1+\cdots+t^DB_D\), we write \(A\degen B\).
The individual substitutions may be Laurent polynomials; only their
combined tensor is required to have no negative powers.
Every degeneration below is supplied with its substitutions and a degree
bound. No general approximation theorem is needed to use it.

For a positive integer \(s\), the notation \(sT\) means the full direct sum of
\(s\) copies of \(T\), with independent variables on all three legs.
It does not mean scalar multiplication of the coefficients.
Tensor products distribute over these direct sums.
We abbreviate tensor products by juxtaposition when no confusion can arise.
We never cancel a direct-sum multiplicity or a tensor factor from a
restriction.

\begin{lemma}[Composition rules used below]\label{lem:composition}
Suppose a triple of Laurent-polynomial linear substitutions \(\Phi_i(t)\)
transforms a tensor \(A_i\) into
\(G_i(t)=\sum_{d=0}^{D_i}t^dG_{i,d}\).
Let \(P_i\) be a specified restriction with \(P_i(C_i)=A_i\), and
let \(R_i\) be a triple of linear maps on the whole target variable spaces.
If \(P_i,R_i\) are independent of \(t\), then
\[
 R_i\Phi_i(t)P_i(C_i)=\sum_{d=0}^{D_i}t^dR_i(G_{i,d}).
\]
Their tensor product has degree at most \(\sum_iD_i\).
Its constant coefficient is \(\bigotimes_iR_i(G_{i,0})\).
Permuting the three legs preserves these statements.
\end{lemma}
\begin{proof}
The first identity follows by substitution and linearity on the tensor
space. In a tensor product of polynomials with nonnegative exponents,
degrees add, and a term has degree zero exactly when every factor does.
Permuting the legs only permutes the three substitution maps.
\end{proof}

The square matrix multiplication tensor is
\[
 \MM_d=\sum_{i,j,k=0}^{d-1}x_{ij}y_{jk}z_{ki}.
\]
Its third variables encode the output coordinates by duality.
Let \(B_X(d)=x\sum_{i=0}^{d-1}y_i z_i\), and define \(B_Y(d)\) and
\(B_Z(d)\) by permuting the legs. Relabeling coordinates gives
\begin{equation}\label{eq:matrix-products}
 \MM_d\MM_e\cong\MM_{de},\qquad
 B_X(d)B_Y(d)B_Z(d)\cong\MM_d.
\end{equation}
For example, the three paired indices in the second product can be taken
as the row, inner, and column indices of matrix multiplication.
Consequently the product of all six leg permutations of \(B_X(d)\)
is \(\MM_{d^2}\).

The basic convolution tensor is
\begin{equation}\label{eq:convolution}
 C(r,b)=\sum_{i=0}^{r-1}\sum_{j=0}^{b-1}x_i y_j z_{i+j}.
\end{equation}
It represents multiplication of polynomials of degrees at most \(r-1\)
and \(b-1\). Put \(d=r+b-1\) and let \(\xi\) be a primitive
\(d\)-th root of unity. Evaluation at \(\xi^\ell\) gives the explicit
rank decomposition
\begin{equation}\label{eq:convolution-rank}
 C(r,b)=\frac1d\sum_{\ell=0}^{d-1}
 \left(\sum_i\xi^{\ell i}x_i\right)
 \left(\sum_j\xi^{\ell j}y_j\right)
 \left(\sum_{h=0}^{d-1}\xi^{-\ell h}z_h\right).
\end{equation}
Indeed \(0\le i+j\le d-1\), so the Fourier sum selects exactly
\(h=i+j\). The output flattening has rank \(d\), and hence
\(\R(C(r,b))=d\).

For \(\pi\in S_3\), superscript \(\pi\) denotes a permutation of the
legs. Define the balanced tensor
\begin{equation}\label{eq:balanced}
 Q_{r,b}=\bigotimes_{\pi\in S_3} C(r,b)^\pi.
\end{equation}
The six factors are kept as labeled factors, even when some are isomorphic.
Swapping the inputs in \eqref{eq:convolution} and reordering these factors
gives
\begin{equation}\label{eq:balanced-properties}
 Q_{r,b}\cong Q_{b,r},\qquad
 \R(Q_{r,b})\le(r+b-1)^6,\qquad Q_{1,2}\cong\MM_4.
\end{equation}
All these isomorphisms are coordinate permutations.

Arithmetic complexity counts scalar additions, subtractions, and
multiplications, including multiplication by fixed coefficients.
The algorithm may depend on a fixed exponent allowance \(\eps\), but
not on the input dimension \(n\). Its constant preprocessing cost may
depend on \(\eps\). This is an arithmetic-operation statement, not
a bit-complexity or floating-point stability statement.

\begin{lemma}[Exactification]\label{lem:exactification}
Suppose specified substitutions transform \(A\) into a polynomial
\(G(t)=\sum_{h=0}^D t^hG_h\). For every positive integer \(k\),
\[
 (kD+1)A^{\otimes k}\restr G_0^{\otimes k}.
\]
\end{lemma}
\begin{proof}
Set \(E=kD+1\) and take a primitive \(E\)-th root \(\eta\).
The tensor polynomial \(G(t)^{\otimes k}\) has degree at most \(kD\).
Thus
\[
 \frac1E\sum_{h=0}^{E-1}G(\eta^h)^{\otimes k}=G_0^{\otimes k}.
\]
Each evaluation uses one copy of \(A^{\otimes k}\); identifying its
output variable groups and absorbing \(1/E\) into one leg gives the
claimed restriction. Every evaluation point is nonzero, so Laurent
substitutions cause no problem.
\end{proof}
